\part{Getting Started}
%   Why LaTeX?
\section{Why \LaTeX?}
% NOTE: Maybe start with something like: "If you ever had to deal with a large Word document, you know it can be a pain to update all its references, replace an image, or just navigate it in general."
% NOTE: LaTeX is one of many tools usable for writing text. 
%       Contrary to other tools such as Word, Google Docs and similar, it is similar to programming.
%       Describe layout with its syntax, write "annotated" text separately with functions -- or as they're called in LaTeX: commands.
%       LaTeX then converts text with its "annotations" together with the layout rules into a PDF (or other output format).
%       Pros: 
%       - Less hassle with updating references
%           - Automatic recomputation of ToC, List of figure, references etc.
%       - Easier structuring for large documents
%       - Uniform appearance throughout document
%       - Easier reference/bibliography management
%       - TODO: add something about handling of Mathematical Formulas
%
%       Cons:
%       - No instant preview -> need to run LaTeX compiler to see result
%       TODO: find more Cons...
%       - Graphical customization can be a pain to do yourself

%   Short history of LaTeX
\section{A Short History of \LaTeX}

%   Setup
\section{Setup}
\subsection{Online -- Overleaf}

\subsection{Windows -- \TeX\ Live}
\subsection{Windows -- MiK\TeX}

\subsection{Linux -- \TeX\ Live}

\subsection{MacOS -- Mac\TeX}

%   Syntax basics
\section{Syntax Basics}
% NOTE: Reserved characters (\, &, #, %, _, ^, {, }, @ etc.)
%       Command dissection (mandatory args, optional args)

%   Minimal Document
\section{Minimal Document}

